\section*{Capítulo \ref{ch:b2:mass-spec-atomic} --- Espectrometría de masas y espectroscopia atómica}

\begin{solution}{exo:b2:mass-spec-atomic:1}
78; $400/2 = 200$; $194 + 1 = 195$, $[\mathrm M + \mathrm H]^+$.
\end{solution}

\begin{solution}{exo:b2:mass-spec-atomic:2}
Masa impar: un número impar de átomos de nitrógeno, uno para una molécula tan pequeña. \ce{C3H7NO} (propanamida, 73)
o \ce{C4H11N} (butan-1-amina, 73).
\end{solution}

\begin{solution}{exo:b2:mass-spec-atomic:3}
Casi iguales: un bromo. 3 : 1: un cloro.
\end{solution}

\begin{solution}{exo:b2:mass-spec-atomic:4}
Amarillo: sodio, \qty{589.0}{nm} y \qty{589.6}{nm}. Carmesí: litio, \qty{670.8}{nm}.
% ledger: asd:NaD2, asd:NaD1, asd:Li671
\end{solution}

\begin{solution}{exo:b2:mass-spec-atomic:5}
$n_C \approx 6.7/1.11 \approx 6$.
\end{solution}

\begin{solution}{exo:b2:mass-spec-atomic:6}
\ce{CO}: 27.9949; \ce{N2}: 28.0061; \ce{C2H4}: 28.0313 u. Poderes de resolución: \ce{CO}/\ce{N2} $28/0.0112
\approx 2.5 \times 10^3$; \ce{N2}/\ce{C2H4} $28/0.0252 \approx 1.1 \times 10^3$; \ce{CO}/\ce{C2H4}
$28/0.0364 \approx 7.7 \times 10^2$.
% ledger: mass:C12, mass:O16, mass:N14, mass:H1
\end{solution}

\begin{solution}{exo:b2:mass-spec-atomic:7}
86: el \omterm{def:b2:mass-spec-atomic:molecular-ion}{ion molecular} de \ce{C5H10O}. 57: \ce{C2H5CO+}, un \omterm{def:b2:aromatic-substitution:friedel-crafts}{ion acilio} procedente de una ruptura $\alpha$ (pérdida de
un radical etilo, 29). 29: \ce{C2H5+}.
% ledger: ms:pentan-3-one
\end{solution}

\begin{solution}{exo:b2:mass-spec-atomic:8}
86: $\mathrm M^{+\bullet}$. 71: $\mathrm M - \ce{CH3}$, el acilio \ce{C3H7CO+}. 43: \ce{CH3CO+}
(pérdida del radical propilo). 58: \omterm{def:b2:mass-spec-atomic:fragmentation}{transposición de McLafferty}, a través de un hidrógeno del carbono $\gamma$ de la
cadena propilo, con pérdida de eteno. En la pentan-3-ona los grupos etilo no tienen carbono $\gamma$: no hay
ion de McLafferty; su pequeño pico en 58 es el pico isotópico \ce{^{13}C} del ion en 57 (tres carbonos,
alrededor del 3.3~\%).
% ledger: ms:pentan-2-one, ms:pentan-3-one
\end{solution}

\begin{solution}{exo:b2:mass-spec-atomic:9}
Una recta que pasa por el origen: pendiente $\sum c_iA_i/\sum c_i^2 = 1.553/30.0 \approx \qty{0.0518}{L/mg}$. Muestra:
$0.128/0.0518 \approx \qty{2.47}{mg/L}$ de cobre.
\end{solution}

\begin{solution}{exo:b2:mass-spec-atomic:10}
$t = L\sqrt{m/(2eU)} = 1.00\sqrt{1000 \times 1.6605 \times 10^{-27}/(2 \times 1.6022 \times 10^{-19}
\times 2.00 \times 10^4)} \approx \qty{16.1}{\micro s}$. Para 1001 u, $t$ crece en el factor
$\sqrt{1.001} \approx 1 + 0.0005$: $\Delta t \approx \qty{8.0}{ns}$.
% ledger: const:u, const:e
\end{solution}

\begin{solution}{exo:b2:mass-spec-atomic:11}
Dos cloros: $(0.758 + 0.242x)^2$ da $0.575$, $0.367$, $0.0586$: $m/z$ 84, 86 y 88 en la razón
100 : 64 : 10.
\end{solution}

\begin{solution}{exo:b2:mass-spec-atomic:12}
\ce{^{40}Ar^{16}O+}: $39.9624 + 15.9949 = 55.9573$; \ce{^{56}Fe}: 55.9349; $R = 56/0.0224 \approx 2.5 \times
10^3$. Si no, mídase otro isótopo del hierro, el \ce{^{57}Fe}, o elimínese el \ce{ArO+} en una celda de colisión
antes del analizador.
% ledger: mass:Ar40, mass:Fe56, mass:O16
\end{solution}

\begin{solution}{pb:b2:mass-spec-atomic:1}
\textbf{1.} 156, 158 y 160 (con 157 y 159): la molécula existe en varias formas isotópicas.
\textbf{2.} Tres picos separados dos unidades, el del medio el mayor: un cloro y un bromo.
\textbf{3.} 156.
\textbf{4.} Masa par: ningún nitrógeno, o un número par.
\textbf{5.} $156 - 79 - 35 = 42$: \ce{C3H6}.
\textbf{6.} \ce{C3H6BrCl}; grado de insaturación $(2 \times 3 + 2 - 6 - 2)/2 = 0$: ni anillo ni doble
enlace.
\textbf{7.} $0.5/14.5 \approx 3.4~\%$, $3.4/1.11 \approx 3$ carbonos.
\textbf{8.} El pico en 157 se da con una precisión de 0.1 en una escala en la que vale 0.5: una incertidumbre del 20~\%.
\textbf{9.} \ce{C3H6^{81}Br^{35}Cl} y \ce{C3H6^{79}Br^{37}Cl}.
\textbf{10.} $3 \times 12 + 6 \times 1.007825 + 78.918338 + 34.968853 \approx 155.9341$.
\textbf{11.} $156/(156.151 - 155.934) \approx 7.2 \times 10^2$.
\textbf{12.} Sí: dos hidrógenos menos, 154.
\textbf{13.} Se pierde un átomo de bromo (79 u 81): \ce{C3H6Cl+}, 77 con \ce{^{35}Cl}, 79 con \ce{^{37}Cl},
3 : 1.
\textbf{14.} Pérdida de HBr: el catión radical $\mathrm{C_3H_5Cl^{+\bullet}}$, de masa par.
\textbf{15.} \ce{C3H5+}, el catión alilo, tras la pérdida de ambos halógenos (Br, y luego HCl).
\textbf{16.} \ce{CH2Cl+} (49 y 51, 3 : 1).
\textbf{17.} \ce{CH2Br+} (93 y 95) y \ce{C2H4Br+} (107 y 109), cada uno 1 : 1: bromo en un extremo de
la cadena; el segundo, por pérdida de \ce{CH2Cl}.
\textbf{18.} El enlace \ce{C-Br} es más débil que el enlace \ce{C-Cl}, y \ce{Br} es el mejor radical
saliente.
\textbf{19.} No: no hay picos en 127--131; \ce{CH2Cl+} y \ce{CH2Br+} muestran cada halógeno en un \ce{CH2}
terminal.
\textbf{20.} \ce{BrCH2CH2CH2Cl}, 1-bromo-3-cloropropano.
\textbf{21.} 156/158/160: M; 77/79: M $-$ Br; 76: M $-$ HBr; 107/109: M $-$ \ce{CH2Cl};
93/95: \ce{CH2Br+}; 49/51: \ce{CH2Cl+}; 41: \ce{C3H5+}; 39, 27: \ce{C3H3+}, \ce{C2H3+}.
\textbf{22.} No: no hay grupo carbonilo.
\textbf{23.} 1-Bromo-2-cloropropano, \ce{CH3CHClCH2Br}: la pérdida de \ce{CH2Br} da \ce{CH3CHCl+} en 63/65.
\textbf{24.} \ce{C3H6^{81}Br^{37}Cl}.
\textbf{25.} $M$: $0.758 \times 0.5069 = 0.384$; $M+2$: $0.758 \times 0.4931 + 0.242 \times 0.5069 =
0.496$; $M+4$: $0.242 \times 0.4931 = 0.119$. Razón $\boldsymbol{\approx 3.2 : 4.2 : 1}$, aproximadamente 3 : 4 :
1; el espectro da $14.5 : 18.6 : 4.4 = 3.3 : 4.2 : 1$.
\end{solution}
